\section{Benchmark design}
\label{sec:design}
\subsection{Towers and extraction points}
Table~\ref{tab:roster} lists the six towers. Each receives an image resized and
center-cropped to $448\times448$, with its native normalization and input packing.
At layer $i$ of a tower with $L$ blocks, relative depth is $i/L$. We sample
$i=\operatorname{round}(Lk/8)$ for $k=1,\ldots,8$; actual fractions vary slightly with $L$.
The deepest extraction includes final tower normalization where present; Qwen3.8
is read at the last block output, before normalization on the merger path.

\begin{table}[t]
\centering\small
\caption{Tower architectures. Tokens are measured before merging at $448\times448$.
The last column is the projector's output width. DINOv2 and SigLIP2 have no native
connector in this study.}
\label{tab:roster}
\begin{tabular}{lrrrrr}
\toprule
Tower & Layers & Width & Patch & Tokens & Projected width\\
\midrule
DINOv2-L & 24 & 1024 & 14 & 1024 & ---\\
SigLIP2-so400m & 27 & 1152 & 14 & 1024 & ---\\
Kimi K2.6 & 27 & 1152 & 14 & 1024 & 7168\\
Kimi K3 & 27 & 1024 & 14 & 1024 & 7168\\
Muse Glimmer & 50 & 1536 & 14 & 1024 & 6656\\
Qwen3.8 & 27 & 1152 & 16 & 784 & 5120\\
\bottomrule
\end{tabular}
\end{table}

For each multimodal tower, we also extract \texttt{merged}, which concatenates each
$2\times2$ group of patch tokens, and \texttt{projected}, the learned map into the
parent language model's embedding width. These stages exist only after the final tower
layer. We do not apply projectors to intermediate layers or exchange them between towers.
Thus the experiment separates extraction points within native pipelines, rather than
independently crossing every tower, layer, and projector.

\subsection{Tasks and readouts}
Recognition uses the 13,000-image validation export of \texttt{ilee0022/ImageNet100},
a 100-class ImageNet subset \citep{deng2009imagenet,ilee0022imagenet100}, split by us into
9,100 training, 1,950 validation, and 1,950 test images. A linear classifier reads either
the mean feature or a learned attention pool. Both operate on a cached $4\times4$ pooled
grid. We report top-1 accuracy after 20 training epochs. Label-budget runs keep validation
and test fixed and train on 100, 455, 1,820, or 9,100 images. The nominal 1\% condition
uses one image per class, hence 1.10\% of the training split.

Geometry uses DIODE depth and normals \citep{vasiljevic2019diode}, with 541 training,
115 validation, and 115 test images drawn from its validation release. The test split
contains 51 indoor and 64 outdoor images. A Probe3D-style multiscale decoder reads the
full patch grid from one layer and trains for 10 epochs. Depth is summarized by
$\delta_1$, the mean per-image fraction of valid pixels satisfying
$\max(\hat d/d,d/\hat d)<1.25$, for prediction $\hat d$ and target $d$; higher is better. Normals use mean angular error in
degrees; lower is better. Targets cover the same crop as the input image.

Correspondence uses cosine feature matching on NAVI, ScanNet, and SPair-71k
\citep{jampani2023navi,dai2017scannet,min2019spair}. NAVI measures recall within 2\,cm
across views of an object; ScanNet measures recall within 10 pixels across indoor views;
SPair measures keypoint accuracy within 10\% of the target bounding-box scale across
object instances. NAVI and ScanNet match on a $64\times64$ feature grid; SPair matches
annotated keypoints against each stage's own grid. These metrics are not averaged together. Occlusion measures top-1 accuracy after masking
0, 10, 20, 35, or 50\% of image area, with readouts trained once on clean images.

KITTI provides a separate depth experiment with 700 training, 150 validation, and
150 test images \citep{geiger2012kitti,uhrig2017sparsity}. We retrain the depth readout
on KITTI to test whether the DIODE stage ordering recurs. This is not zero-shot transfer
of a DIODE-trained readout. Dataset preparation and evaluation details appear in
Appendix~\ref{app:protocol}.

\subsection{Capacity, selection, and uncertainty}
\label{sec:selection}
The core recognition and DIODE sweeps run at native feature width (\emph{raw}) and after
training-split PCA to 512 dimensions (\emph{matched}). PCA is frozen before readout training. Matching
holds head parameter count constant within a readout type; it does not equate spatial
sampling or the information retained by the reduction. We retain both arms as a
sensitivity analysis.

A cell specifies tower, extraction point, task/readout, and capacity arm. The core sweeps
contain 224 recognition and 224 geometry cells. Each selects its learning rate on
validation from a shared task-specific grid, then trains with three seeds. Headline
intervals use 10,000 paired bootstrap resamples at 95\% confidence. For trained readouts,
we first average each image's score over seeds, then resample images; correspondence
resamples image pairs. These intervals condition on the fitted runs and do not include
all training uncertainty. We report seed standard deviations separately. Intervals are
pointwise, without a multiple-comparison correction. An interval crossing zero is
unresolved, not evidence of equivalence; an untested pair has no inferential ranking.

\subsection{Controls and corrections}
\label{sec:controls}
Extraction tests compare weights with parent checkpoints and restore spatial token order
before dense evaluation. Kimi K3 has tower-weight equality but lacks the forward-output
and projector-weight checks available for the other extracted models. Lossless merging
is verified against cached tensors for the two Kimi towers and inferred from architecture
for Muse Glimmer and Qwen3.8.

The pooled recognition cache has an incomplete control. On SigLIP2 and Muse Glimmer,
mean-readout results agree with full-token features, but six of sixteen full-token
attention cells collapse during optimization. Follow-up runs do not establish stable
agreement. Attention-based recognition levels and layer curves therefore remain
conditional on pooled features.

Two corrections changed the recorded results: expanding the recognition learning-rate
grid and seeding randomized PCA. Initially, 107 of 224 cells selected the lower rate
limit. In the corrected run none selects the lower limit and three select the upper
limit. All matched recognition cells and headline geometry comparisons were rerun. Together these corrections changed an apparent Muse
Glimmer merge-to-projector loss from $-0.0075$ to $+0.0042$ top-1. The full matched geometry sweep retains its original
randomized reductions. Its curves and selected peaks are descriptive; geometry intervals
refer to deterministic reruns. Appendix~\ref{app:corrections} records the corrections.
